% Chapter — Axisymmetric ascent aerodynamics (FR-9). Follows the mandatory
% FR-29 six-section template: purpose; assumptions and domain of validity;
% derivation; discretization and implementation; validation evidence;
% references. The equation labels eq:aero:mach, eq:aero:qbar,
% eq:aero:alpha, eq:aero:interp, eq:aero:axial, eq:aero:normal,
% eq:aero:cpmoment, and eq:aero:damping are echoed verbatim in comments in
% cpp/src/models/aero.cpp (FR-29 equation-label-to-code traceability);
% renaming them breaks that trace.
\chapter{Axisymmetric Ascent Aerodynamics}
\label{ch:aero}

\section{Purpose}
\label{sec:aero:purpose}

This chapter documents the FR-9 continuum ascent aerodynamics model: the
body-frame aerodynamic force and the moment about the current composite
center of gravity (CG) of an axisymmetric launch vehicle, evaluated from a
per-stack-configuration Mach-table database $C_A(M)$, $C_{N\alpha}(M)$,
$x_{cp}(M)$ with an optional constant pitch-damping derivative $C_{mq}$,
using the total-angle-of-attack formulation. The database is configured by
the FR-13 vehicle schema's \texttt{[[aero]]} blocks and their CSV Mach
tables; this model is the ascent-regime counterpart of the cannonball
orbital-drag model (Chapter~\ref{ch:drag}), and the two regimes never
apply simultaneously (the FR-9 domain split). The model is implemented in
\texttt{cpp/src/models/aero.cpp} with its interface in
\texttt{cpp/include/star/models/aero.hpp}.

\section{Assumptions and domain of validity}
\label{sec:aero:domain}

Assumptions:
\begin{enumerate}
  \item \textbf{Axisymmetric configuration, total-angle-of-attack
    loading.} The stack is aerodynamically axisymmetric about the
    structural $+X$ axis (FR-13: $+X$ toward the nose, origin at the aft
    plane). The loading then depends on Mach number and on the single
    total angle of attack $\alpha_t$ between $+\hat{\vv{x}}$ and the
    air-relative velocity; by symmetry the resultant force lies in the
    wind plane spanned by $\hat{\vv{x}}$ and $\vv{v}_{\mathrm{rel}}$.
    There is no side force out of that plane and no aerodynamic roll
    torque (no fins, no $C_{lp}$): the torque $x$-component is
    structurally zero.
  \item \textbf{Air-relative velocity supplied by the caller (the FR-8
    rule).} All aerodynamic quantities --- Mach number, dynamic pressure,
    angle of attack --- use the air-relative velocity of a co-rotating
    atmosphere, $\vv{v}_{\mathrm{rel}} = \vv{v} - \vv{\omega}_p \times
    \vv{r}$ (Chapter~\ref{ch:drag}, equation eq:drag:vrel), which the
    equations-of-motion layer computes and rotates into the body frame;
    this module is frame-agnostic given the body-frame
    $\vv{v}_{\mathrm{rel}}$. Winds are neglected, as in
    Chapter~\ref{ch:drag}. The density $\rho$ and speed of sound $a$ come
    from the atmosphere models (Chapter~\ref{ch:ussa76} provides both for
    Earth ascent).
  \item \textbf{Normal force linear in $\sin\alpha_t$ (documented
    choice).} The database carries the normal-force slope
    $C_{N\alpha}(M)$ per radian and the model evaluates
    $C_N = C_{N\alpha}(M)\,\sin\alpha_t$. This is the standard
    slender-body crossflow form: it is odd in $\alpha_t$, matches the
    linear-in-$\alpha$ slope at small angles to within
    $\alpha_t^3/6$ (relative error $\le 0.5\,\%$ at
    $\alpha_t = 10^\circ$), is algebraically free of normalization
    singularities (Section~\ref{sec:aero:derivation}), and degrades
    gracefully toward the crossflow-only geometry at
    $\alpha_t \to 90^\circ$. The database itself is fitted for the
    small-to-moderate angles of ascent flight; evaluations at large
    $\alpha_t$ are a geometric extension of the formulation, not data.
  \item \textbf{Composite CG on the symmetry axis.} The moment uses only
    the axial CG station $x_{cg}$: the axial force acts along the
    symmetry axis and therefore produces no moment about an on-axis CG,
    and the normal force acts at the center-of-pressure station
    $x_{cp}(M)$ on the axis. Lateral CG offsets of the FR-13 stacks
    (whose sub-blocks are nominally axisymmetric) are neglected; the
    induced moment error is of order (lateral offset) $\times$ (force),
    second-order for near-axis CGs.
  \item \textbf{Quasi-steady coefficients with optional constant pitch
    damping.} Unsteady effects are collapsed into the single optional
    constant $C_{mq}$ (per radian of nondimensional pitch rate
    $\hat{q} = \omega\, d_{\mathrm{ref}} / (2\lVert \vv{v}_{\mathrm{rel}}
    \rVert)$, referenced to $S_{\mathrm{ref}}$ and the reference diameter
    $d_{\mathrm{ref}}$ --- the convention declared for the vehicle
    schema's \texttt{cmq\_per\_rad} key). $C_{mq} \le 0$ opposes the
    transverse rate; exactly zero disables the term. The roll rate is
    undamped (assumption 1).
  \item \textbf{Piecewise-linear Mach interpolation, clamped ends
    (documented approximation).} Coefficients between breakpoints are
    piecewise-linear in Mach; below the first and above the last
    breakpoint the end row is held constant. Both choices are
    deterministic, continuous, and exact at the breakpoints
    (Section~\ref{sec:aero:implementation}).
\end{enumerate}

Domain of validity: continuum ascent flight only --- the orbital
free-molecular regime is Chapter~\ref{ch:drag}'s cannonball model (FR-9
domain split) --- in forward flight at the small-to-moderate total angles
of attack for which a linear normal-force-slope database is meaningful,
with Mach inside the configured table (clamped outside, per assumption 6).
Zero air-relative velocity (the vehicle on the pad before release) returns
exactly zero force, moment, Mach, dynamic pressure, and angle of attack.
Out-of-domain behavior, matching the code exactly: malformed tables
(fewer than two rows, column-length mismatches, a non-increasing or
negative Mach grid, non-finite entries, non-positive $S_{\mathrm{ref}}$ or
$d_{\mathrm{ref}}$, positive $C_{mq}$) and malformed inputs (non-finite
velocity, rates, or CG station; negative or non-finite $\rho$;
non-positive or non-finite $a$; a negative or non-finite Mach passed to
the coefficient lookup) throw \texttt{std::domain\_error}.

\section{Derivation}
\label{sec:aero:derivation}

\subsection{Air data and total angle of attack}

With the body-frame air-relative velocity $\vv{v}_{\mathrm{rel}} =
(v_x, v_y, v_z)^{\mathsf T}$, speed
$v = \lVert \vv{v}_{\mathrm{rel}} \rVert$, and the caller-supplied speed
of sound $a$ and density $\rho$:
\begin{equation}
  M \;=\; \frac{v}{a},
  \label{eq:aero:mach}
\end{equation}
\begin{equation}
  \bar{q} \;=\; \tfrac{1}{2}\, \rho\, v^2 .
  \label{eq:aero:qbar}
\end{equation}
Decompose $\vv{v}_{\mathrm{rel}}$ into its axial part $v_x \hat{\vv{x}}$
and the crossflow part
$\vv{v}_\perp = (0, v_y, v_z)^{\mathsf T}$. The total angle of attack is
the angle between the symmetry axis and the wind,
\begin{equation}
  \alpha_t \;=\; \operatorname{atan2}\!\bigl(
    \lVert \vv{v}_\perp \rVert,\; v_x \bigr)
  \;\in\; [0, \pi],
  \qquad
  \sin\alpha_t = \frac{\lVert \vv{v}_\perp \rVert}{v}, \quad
  \cos\alpha_t = \frac{v_x}{v}.
  \label{eq:aero:alpha}
\end{equation}

\subsection{Force decomposition}

For an axisymmetric body the pressure and shear distribution is invariant
under a rotation about $\hat{\vv{x}}$ applied jointly to the body and the
wind vector, so the resultant force lies in the wind plane and is fully
described by two dimensionless coefficients on the dynamic pressure and
the configured reference area $S_{\mathrm{ref}}$
\cite[Sect.~3.5]{montenbruck2000}, \cite[Sect.~8.6.2]{vallado2013}: an
axial force coefficient $C_A$ along $-\hat{\vv{x}}$ and a normal force
coefficient $C_N$ in the wind plane, perpendicular to the axis, opposing
the crossflow. With the database convention of assumption~3,
$C_N = C_{N\alpha}(M) \sin\alpha_t$:
\begin{equation}
  \boxed{\;
  \vv{F}_A \;=\; -\,\bar{q}\, S_{\mathrm{ref}}\, C_A(M)\, \hat{\vv{x}},
  \;}
  \label{eq:aero:axial}
\end{equation}
\begin{equation}
  \boxed{\;
  \vv{F}_N \;=\; -\,\bar{q}\, S_{\mathrm{ref}}\, C_{N\alpha}(M)\,
    \sin\alpha_t\; \hat{\vv{v}}_\perp
  \;=\; -\,\frac{\bar{q}\, S_{\mathrm{ref}}\, C_{N\alpha}(M)}{v}\,
    \vv{v}_\perp ,
  \;}
  \label{eq:aero:normal}
\end{equation}
where the second form of \eqref{eq:aero:normal} substitutes
$\sin\alpha_t$ from \eqref{eq:aero:alpha} and absorbs the crossflow
normalization: it involves no trigonometric evaluation and no division by
$\lVert \vv{v}_\perp \rVert$, so zero crossflow yields an exactly zero
normal force with no special case, and the only excluded point of the
whole formulation is $v = 0$, which returns the structural zero of
Section~\ref{sec:aero:domain}. The total force is
$\vv{F} = \vv{F}_A + \vv{F}_N$.

\subsection{Center-of-pressure moment}

The database locates the resultant normal force at the
center-of-pressure station $x_{cp}(M)$ on the symmetry axis, expressed in
the same structural frame as every vehicle station (origin at the aft
plane, $+X$ toward the nose --- exactly the \texttt{xcp\_m} column the
vehicle schema declares). With the composite CG at station $x_{cg}$ on
the axis (assumption 4), the moment about the CG is the lever arm crossed
with the normal force:
\begin{equation}
  \boxed{\;
  \vv{m}_{cp} \;=\; \bigl( x_{cp}(M) - x_{cg} \bigr)\, \hat{\vv{x}}
    \times \vv{F}_N .
  \;}
  \label{eq:aero:cpmoment}
\end{equation}
The axial force contributes no moment (it acts along the axis through the
on-axis CG), so $\alpha_t = 0$ produces exactly zero static moment. When
the center of pressure is forward of the CG ($x_{cp} > x_{cg}$),
\eqref{eq:aero:cpmoment} is an overturning (destabilizing) moment --- the
normal case for the finless launchers of the starter fleet, which are
stabilized by thrust-vector control; $x_{cp} < x_{cg}$ gives a restoring
moment.

\subsection{Pitch-damping moment}

The optional damping term models the moment that opposes the transverse
angular rate. Using the classical rate nondimensionalization
$\hat{q} = \omega\, d_{\mathrm{ref}} / (2 v)$ with the configured
reference diameter $d_{\mathrm{ref}}$, the damping contribution to the
moment coefficient about each transverse axis is
$C_{mq}\,\hat{q}$, i.e.\ a dimensional moment
$\bar{q} S_{\mathrm{ref}} d_{\mathrm{ref}} \cdot C_{mq}\,
\omega d_{\mathrm{ref}}/(2v)$ per transverse axis. Applied isotropically
in the transverse plane (axisymmetry), with
$\vv{\omega}_\perp = (0, \omega_y, \omega_z)^{\mathsf T}$ the transverse
part of the body angular velocity:
\begin{equation}
  \boxed{\;
  \vv{m}_q \;=\; \frac{\bar{q}\, S_{\mathrm{ref}}\,
    d_{\mathrm{ref}}^2\, C_{mq}}{2\, v}\; \vv{\omega}_\perp .
  \;}
  \label{eq:aero:damping}
\end{equation}
$C_{mq} \le 0$ makes \eqref{eq:aero:damping} dissipative; $C_{mq} = 0$
disables the term exactly, and a pure roll rate
($\vv{\omega}_\perp = \vv{0}$) produces exactly zero damping moment. The
total moment about the CG is $\vv{m} = \vv{m}_{cp} + \vv{m}_q$.

\subsection{Mach-table interpolation}

Each coefficient column $c \in \{C_A, C_{N\alpha}, x_{cp}\}$ is tabulated
on one strictly increasing Mach grid $M_1 < M_2 < \dots < M_n$
($n \ge 2$, $M_1 \ge 0$). The lookup is piecewise linear with clamped
ends:
\begin{equation}
  c(M) \;=\;
  \begin{cases}
    c_1, & M \le M_1, \\[2pt]
    c_i + u\, (c_{i+1} - c_i),
      \quad u = \dfrac{M - M_i}{M_{i+1} - M_i},
      & M_i \le M < M_{i+1}, \\[8pt]
    c_n, & M \ge M_n .
  \end{cases}
  \label{eq:aero:interp}
\end{equation}
The interpolant is continuous everywhere, including at the clamp
junctions, and exact at every breakpoint: the clamp branches return the
end rows literally, and an interior breakpoint selects its own segment
with $u = 0$, so $c_i + u(c_{i+1} - c_i)$ reproduces $c_i$ exactly in
IEEE arithmetic.

\section{Discretization and implementation}
\label{sec:aero:implementation}

\begin{enumerate}
  \item \textbf{Module and traceability.} The model lives in
    \texttt{cpp/src/models/aero.cpp}; the source comments echo the labels
    \texttt{eq:aero:mach} through \texttt{eq:aero:damping} verbatim.
  \item \textbf{Structural zeros.} $v = 0$ returns literal zeros in every
    output before any table access --- the logged dynamic pressure on the
    pad is exactly zero (Phase~4 exit criterion 10). Zero crossflow
    ($v_y = v_z = 0$) skips the normal-force/center-of-pressure block
    entirely, leaving literal zeros (axial force only); $C_{mq} = 0$ or a
    purely axial rate skips the damping block; the torque $x$-component
    is never written (no roll torque, assumption 1). These are branch
    decisions on exact binary64 comparisons, not epsilon tests.
  \item \textbf{Fixed evaluation order.} Segment selection scans the grid
    in ascending order; every product chain is written in one fixed
    association. Together with the exact-breakpoint property of
    \eqref{eq:aero:interp}, evaluation at a committed table breakpoint
    reads the table row exactly.
  \item \textbf{Determinism.} The force/moment path uses IEEE-754 basic
    operations and square roots only (Eigen's \texttt{norm()}), so it is
    bit-portable across platforms under the D-10 flags. The reported
    $\alpha_t$ diagnostic \eqref{eq:aero:alpha} is the module's only libm
    call (\texttt{atan2}); it is never consumed by the force/moment path,
    so its one-ulp cross-platform spread never feeds back into the
    dynamics.
  \item \textbf{Pure function over plain structs.} The table struct
    mirrors the FR-13 vocabulary (\texttt{ref\_area\_m2},
    \texttt{ref\_diameter\_m}, \texttt{cmq\_per\_rad}, and the CSV
    columns \texttt{mach}, \texttt{ca}, \texttt{cnalpha\_per\_rad},
    \texttt{xcp\_m}); the Python layer validates and fills it across the
    binding (D-2: the core never parses text). The caller supplies the
    body-frame $\vv{v}_{\mathrm{rel}}$ (FR-8), $\rho$ and $a$ from the
    atmosphere model, the CG station $x_{cg}$ from the mass-property
    model (Chapter~\ref{ch:massprops}), and the body rates; selecting the
    active \texttt{[[aero]]} configuration across staging events is the
    vehicle layer's responsibility.
\end{enumerate}

\section{Validation evidence}
\label{sec:aero:validation}

Table~\ref{tab:aero:validation} lists the tests that constitute this
chapter's validation evidence. Each test identifier is machine-checked to
exist in the test tree by \texttt{scripts/lint\_chapters.py}; golden
provenance and tolerances are recorded in
\texttt{tests/golden/aero/manifest.toml}.

\begin{table}[htbp]
  \centering
  \caption{Validation evidence for the axisymmetric aero model (doctest,
    \texttt{cpp/tests/test\_aero.cpp}).}
  \label{tab:aero:validation}
  \begin{tabular}{lp{8.6cm}}
    \toprule
    Test ID & Acceptance basis \\
    \midrule
    \testid{AERO-INTERP-GOLDEN} &
      The lookup \eqref{eq:aero:interp} reproduces every committed table
      row exactly (bit equality) at every Mach breakpoint of both fleet
      tables and at the clamped ends, and matches the independent
      60-digit mpmath references to $10^{-12}$ relative at every segment
      midpoint. \\
    \testid{AERO-BREAKPOINT-GOLDEN} &
      Force/moment reconstruction at every
      $C_A/C_{N\alpha}/x_{cp}$ Mach breakpoint of both committed fleet
      tables matches the independent 60-digit mpmath references to
      $10^{-12}$ (norm-)relative on force and moment, and to $10^{-12}$
      relative on Mach, dynamic pressure, and total angle of attack
      (Phase~4 exit criterion 8). \\
    \testid{AERO-FORCETORQUE-GOLDEN} &
      The formulation families --- total-alpha sweep ($0$, $0.5^\circ$,
      $10^\circ$, $90^\circ$, retrograde), crossflow roll orientations,
      CG fore and aft of the center of pressure, damping on/off/roll-only
      per \eqref{eq:aero:damping}, reference-area/diameter scaling,
      off-breakpoint interpolation, and above-table clamping --- match
      the mpmath references to $10^{-12}$, with every reference recorded
      as zero required exactly zero. \\
    \testid{AERO-STRUCTURAL-ZEROS} &
      Zero air-relative velocity returns exactly zero in every output
      including the dynamic pressure (Phase~4 exit criterion 10 hook);
      zero total alpha yields exactly zero normal force and static
      moment; $C_{mq} = 0$ and pure-roll-rate evaluations are
      bit-identical to their damping-free counterparts; the torque
      $x$-component is exactly zero in a generic evaluation. \\
    \testid{AERO-DOMAIN} &
      Out-of-domain inputs per Section~\ref{sec:aero:domain} throw
      \texttt{std::domain\_error}. \\
    \bottomrule
  \end{tabular}
\end{table}

\section{References}
\label{sec:aero:references}

The dynamic-pressure dimensionless-coefficient force convention follows
Montenbruck and Gill~\cite[Sect.~3.5]{montenbruck2000} and
Vallado~\cite[Sect.~8.6.2]{vallado2013}, the same convention already
relied on by Chapter~\ref{ch:drag}. The axisymmetric total-angle-of-attack
decomposition \eqref{eq:aero:axial}--\eqref{eq:aero:normal} is derived in
Section~\ref{sec:aero:derivation} from the symmetry of the configured
vehicle; the center-of-pressure moment \eqref{eq:aero:cpmoment} and the
pitch-damping nondimensionalization \eqref{eq:aero:damping} are the
coefficient-database conventions this project declares for the FR-13
vehicle schema (documented choices of this chapter, not external physical
results). The committed fleet tables carry provenance
``representative'' in their CSV headers. Full bibliographic details
appear in the bibliography.
